%%%%%%%%%%%%%%%%%%%%%%%%%%%%%%%%%%%%%%%%%%%%%%%%%%%%%%%%%%%%%%%%%%%%%
%%                                                                 %%
%% Please do not use \input{...} to include other tex files.       %%
%% Submit your LaTeX manuscript as one .tex document.              %%
%%                                                                 %%
%% All additional figures and files should be attached             %%
%% separately and not embedded in the \TeX\ document itself.       %%
%%                                                                 %%
%%%%%%%%%%%%%%%%%%%%%%%%%%%%%%%%%%%%%%%%%%%%%%%%%%%%%%%%%%%%%%%%%%%%%

%%\documentclass[referee,sn-basic]{sn-jnl}% referee option is meant for double line spacing

%%=======================================================%%
%% to print line numbers in the margin use lineno option %%
%%=======================================================%%

%%\documentclass[lineno,sn-basic]{sn-jnl}% Basic Springer Nature Reference Style/Chemistry Reference Style

%%======================================================%%
%% to compile with pdflatex/xelatex use pdflatex option %%
%%======================================================%%

%%\documentclass[pdflatex,sn-basic]{sn-jnl}% Basic Springer Nature Reference Style/Chemistry Reference Style

%%\documentclass[sn-basic]{sn-jnl}% Basic Springer Nature Reference Style/Chemistry Reference Style
\documentclass[sn-mathphys]{sn-jnl}% Math and Physical Sciences Reference Style
%%\documentclass[sn-aps]{sn-jnl}% American Physical Society (APS) Reference Style
%%\documentclass[sn-vancouver]{sn-jnl}% Vancouver Reference Style
%%\documentclass[sn-apa]{sn-jnl}% APA Reference Style
%%\documentclass[sn-chicago]{sn-jnl}% Chicago-based Humanities Reference Style
%%\documentclass[sn-standardnature]{sn-jnl}% Standard Nature Portfolio Reference Style
%%\documentclass[default]{sn-jnl}% Default
%%\documentclass[default,iicol]{sn-jnl}% Default with double column layout

%%%% Standard Packages
%%<additional latex packages if required can be included here>
%%%%

%%%%%=============================================================================%%%%
%%%%  Remarks: This template is provided to aid authors with the preparation
%%%%  of original research articles intended for submission to journals published
%%%%  by Springer Nature. The guidance has been prepared in partnership with
%%%%  production teams to conform to Springer Nature technical requirements.
%%%%  Editorial and presentation requirements differ among journal portfolios and
%%%%  research disciplines. You may find sections in this template are irrelevant
%%%%  to your work and are empowered to omit any such section if allowed by the
%%%%  journal you intend to submit to. The submission guidelines and policies
%%%%  of the journal take precedence. A detailed User Manual is available in the
%%%%  template package for technical guidance.
%%%%%=============================================================================%%%%

\usepackage[T1]{fontenc}
%%\unnumbered% uncomment this for unnumbered level heads
\usepackage{graphicx}

\usepackage{hyperref} 
\jyear{2021}%

%% as per the requirement new theorem styles can be included as shown below
\theoremstyle{thmstyleone}%
\newtheorem{theorem}{Theorem}%  meant for continuous numbers
%%\newtheorem{theorem}{Theorem}[section]% meant for sectionwise numbers
%% optional argument [theorem] produces theorem numbering sequence instead of independent numbers for Proposition
\newtheorem{proposition}[theorem]{Proposition}%
%%\newtheorem{proposition}{Proposition}% to get separate numbers for theorem and proposition etc.

\theoremstyle{thmstyletwo}%
\newtheorem{example}{Example}%
\newtheorem{remark}{Remark}%

\theoremstyle{thmstylethree}%
\newtheorem{definition}{Definition}%

\raggedbottom


\input{sec/preamble}
\input{sec/notations}

\begin{document}

\title[Article Title]{Article Title}

%%=============================================================%%
%% Prefix   -> \pfx{Dr}
%% GivenName    -> \fnm{Joergen W.}
%% Particle -> \spfx{van der} -> surname prefix
%% FamilyName   -> \sur{Ploeg}
%% Suffix   -> \sfx{IV}
%% NatureName   -> \tanm{Poet Laureate} -> Title after name
%% Degrees  -> \dgr{MSc, PhD}
%% \author*[1,2]{\pfx{Dr} \fnm{Joergen W.} \spfx{van der} \sur{Ploeg} \sfx{IV} \tanm{Poet Laureate}
%%                 \dgr{MSc, PhD}}\email{iauthor@gmail.com}
%%=============================================================%%

\author[1,2]{\fnm{First} \sur{Author}}

\author[2,3]{\fnm{Second} \sur{Author}}
%\equalcont{These authors contributed equally to this work.}

\author[1,2]{\fnm{Third} \sur{Author}}
%\equalcont{These authors contributed equally to this work.}

\affil[1]{\orgdiv{Department}, \orgname{Organization}, \orgaddress{\city{City} \postcode{100190},  \country{Country}}}

\affil[2]{\orgdiv{Department}, \orgname{Organization}, \orgaddress{\city{City} \postcode{10587},  \country{Country}}}

\affil[3]{\orgdiv{Department}, \orgname{Organization}, \orgaddress{\city{City} \postcode{610101},  \country{Country}}}

%%==================================%%
%% sample for unstructured abstract %%
%%==================================%%
\input{sec/0_exp_info}
\input{sec/0_abstract_mir}

%%================================%%
%% Sample for structured abstract %%
%%================================%%

% \abstract{\textbf{Purpose:} The abstract serves both as a general introduction to the topic and as a brief, non-technical summary of the main results and their implications. The abstract must not include subheadings (unless expressly permitted in the journal's Instructions to Authors), equations or citations. As a guide the abstract should not exceed 200 words. Most journals do not set a hard limit however authors are advised to check the author instructions for the journal they are submitting to.
%
% \textbf{Methods:} The abstract serves both as a general introduction to the topic and as a brief, non-technical summary of the main results and their implications. The abstract must not include subheadings (unless expressly permitted in the journal's Instructions to Authors), equations or citations. As a guide the abstract should not exceed 200 words. Most journals do not set a hard limit however authors are advised to check the author instructions for the journal they are submitting to.
%
% \textbf{Results:} The abstract serves both as a general introduction to the topic and as a brief, non-technical summary of the main results and their implications. The abstract must not include subheadings (unless expressly permitted in the journal's Instructions to Authors), equations or citations. As a guide the abstract should not exceed 200 words. Most journals do not set a hard limit however authors are advised to check the author instructions for the journal they are submitting to.
%
% \textbf{Conclusion:} The abstract serves both as a general introduction to the topic and as a brief, non-technical summary of the main results and their implications. The abstract must not include subheadings (unless expressly permitted in the journal's Instructions to Authors), equations or citations. As a guide the abstract should not exceed 200 words. Most journals do not set a hard limit however authors are advised to check the author instructions for the journal they are submitting to.}

\keywords{Keyword1, keyword2, keyword3, keyword4, keywro5}

\maketitle

\input{sec/1_intro_v2}
\input{sec/2_related_work_v1_cat}
% \input{sec/3_methodology}
\input{sec/3_method_v1_cat}
\input{sec/4_experiments_v1_cat}
\section{Conclusion}

In this paper, we have demonstrated that the intrinsic similarity among character components can be leveraged to substantially enhance recognition performance for characters situated in the tail of long-tail distributions. This improvement is especially critical for new or infrequently encountered characters, as their representation is often underemphasized in traditional deep learning models. By recognizing and exploiting these shared structural elements, the network gains an enriched understanding of how individual strokes or radical components form complete characters, ultimately boosting accuracy where it is most needed.

Moreover, we introduced the spindle-shaped network architecture as a powerful means of extracting and refining local character component features. This design strategically expands channel capacity in mid-level layers, facilitating a richer embedding space and promoting more reliable feature representations that capture subtle distinctions between visually similar glyphs. The resultant focus on local components not only bolsters recognition of the rarest characters but also provides a more robust framework for accommodating newly encountered symbols. Through additional analyses, we verified the effectiveness of SpindleNet in adapting to both head and tail classes, ensuring balanced performance across the entire character set.

Extensive experiments conducted on three challenging Chinese ancient book datasets, namely TKH, MTH1000, and MTH1200, further validate the advantages of our method. The outcomes confirm that SpindleNet achieves state-of-the-art performance, surpassing conventional solutions by handling extreme data imbalance and preserving fidelity for high-frequency characters. In challenging historical textual corpora that feature a broad and often sparsely sampled set of characters, such improved recognition capabilities are indispensable for accurate transcription and subsequent scholarly research.

Looking ahead, we plan to extend our framework to a weakly-supervised training paradigm, a direction that has already shown promise in scene text recognition tasks. This approach could reduce the burden of precise annotations while potentially increasing the volume of available training samples. By integrating weak labels or partial annotations, future models may become even more adept at unraveling the intricate forms of rare and unseen characters, thereby advancing the state of the art for ancient document understanding.

\input{sec/6_acknowledgement}

\bibliography{ref/refs,ref/all.bib}

\end{document}
